\chapter{Driver Control}
\label{ch:driver}

In the driver control period, a person steers the robot with the controller.
The left stick drives, and the right stick moves the arm and claw. The program's
job is to turn stick positions into motor power in a way that feels smooth and
predictable. This chapter follows a stick reading through each processing step,
then explains the arm, the claw and the cone assist.

All of this happens inside \code{user_control()} in \code{src/drivercontrol.py},
which repeats every 20~ms (\code{LOOP_MS}), or 50 times per second.

\section{From Stick to Target Power}
\label{sec:stick}

A stick reports its position as a number from $-100$ to $100$. Three steps turn
this into a \emph{target} power for the drive.

\textbf{Deadband.} A stick rarely rests at exactly zero, so a released stick
might report 1 or $-1$. If this were passed to the motors, the robot would creep.
Readings smaller than \code{DEADBAND} (2) are therefore treated as zero.

\textbf{Squaring.} The reading is squared while keeping its sign:
\begin{equation}
  f = \frac{x\,|x|}{100} \times m ,
  \label{eq:square}
\end{equation}
where $x$ is the stick reading and $m$ is the speed multiplier. Squaring gives
the driver finer control near the centre, where precise movements happen, while
still reaching full power at the edge.

\textbf{Worked example:} With the stick halfway forward ($x = 50$) and the normal
multiplier $m = 0.5$, $f = 50 \times 50 / 100 \times 0.5 = 12.5\%$. Without
squaring, the same stick position would give 25\%. So the first half of the
stick's travel covers only the slowest quarter of the speed range, which makes
small adjustments much easier.

\textbf{Speed modes.} The multiplier $m$ depends on the right shoulder buttons:

\begin{table}[h]
  \centering
  \begin{tabular}{@{}lll@{}}
    \toprule
    Held & Multiplier & Use \\
    \midrule
    nothing & 0.50 (\code{SPEED_NORMAL}) & general driving \\
    R1      & 1.00 (\code{SPEED_BOOST})  & crossing the field quickly \\
    R2      & 0.25 (\code{SPEED_CRAWL})  & lining up precisely \\
    \bottomrule
  \end{tabular}
  \caption{Speed modes.}
  \label{tab:speeds}
\end{table}

The turning stick is processed in the same way, except that it is divided by 200
rather than 100, and multiplied by \code{TURN_NORMAL} (0.8) unless R1 is held.
This keeps turns gentler than forward movement in normal driving.

\section{Slew Limiting}
\label{sec:slew}

If the driver pushes the stick from zero to full in an instant, sending full
power straight to the motors would jolt the robot. This can make it tip, spin
its wheels, or throw a game object out of the claw. \emph{Slew limiting} stops
this by letting the power change by only a small step each loop.

\lstinputlisting[firstline=10,firstnumber=10,lastline=22,caption={\code{slew()} in \code{src/drivercontrol.py}, lines 10--22.},label=lst:slew]{\src{drivercontrol.py}}

There are two step sizes. Speeding up uses the smaller \emph{acceleration}
step, so the robot gets moving smoothly. Slowing down (or reversing) uses the
larger \emph{deceleration} step, so the robot still stops quickly when the
driver lets go. Line 12 decides which one applies: if the target has the
opposite sign to the current power, or is smaller in size, the robot is slowing
down.

\textbf{Worked example:} \code{FWD_ACCEL_STEP} is 4, so going from 0 to 100\% takes
$100 / 4 = 25$ loops, or $25 \times 20~\text{ms} = 500$~ms.
\code{FWD_DECEL_STEP} is 12, so stopping from 100\% takes about 9 loops, or
170~ms. Turning uses larger steps (8 and 20), because a slow-to-respond turn
feels sluggish to the driver.

\section{Balancing Forward and Turn}
\label{sec:taper}

In arcade drive, the left motor gets \textit{forward + turn} and the right motor
gets \textit{forward $-$ turn}. Two problems come from this.

\textbf{Turning too sharply at speed.} The same turn stick input that is useful
for spinning on the spot makes the robot swerve violently at full speed.
\code{turn_taper()} reduces the turn as the forward speed increases:
\begin{equation}
  g(f) = g_{\text{rest}} + (g_{\text{full}} - g_{\text{rest}}) \frac{|f|}{100} ,
  \label{eq:taper}
\end{equation}
with $g_{\text{rest}} = 1.00$ and $g_{\text{full}} = 0.35$. At rest the driver
gets the full turn; at full speed only 35\% of it; and at 50\% forward power,
$g = 1 + (0.35 - 1)(0.5) = 0.675$.

\textbf{Clipping.} A motor cannot exceed 100\%. If forward is 80 and turn is 40,
the left motor would be asked for 120 and is cut to 100, while the right motor
gets 40. The robot's average forward speed drops from 80 to 70, so the driver
loses speed just by turning. \code{fit_turn()} prevents this by shrinking the turn
instead of the forward power:

\lstinputlisting[firstline=37,firstnumber=37,lastline=44,caption={\code{fit_turn()} in \code{src/drivercontrol.py}, lines 37--44.},label=lst:fitturn]{\src{drivercontrol.py}}

\textbf{Worked example:} With forward 80 and turn 40, there is $100 - 80 = 20$
of room left. With \code{FORWARD_PRIORITY} at 0.8, the allowed turn is
$40 + (20 - 40)(0.8) = 24$. The motors receive $80 + 24 = 104$ (cut to 100) and
$80 - 24 = 56$, an average of 78 instead of 70. Most of the forward speed is
kept, and the robot still turns.

Holding R1 skips both the taper and \code{fit_turn()}, giving the driver raw,
unrestricted control when they want it.

\section{Arm and Claw}
\label{sec:manipulator}

The right stick's vertical axis moves the arm and its horizontal axis moves the
claw. The arm is squared and slew-limited in the same way as the drive, because
a heavy arm that starts or stops suddenly swings and can shake loose what it is
carrying. The claw is not slew-limited, since it is light and needs to respond
immediately.

Both use \code{HOLD} when stopped (\code{manipulator_init()} in
\code{config.py}), so the arm does not sag under its own weight. The claw's
torque is limited with \code{CLAW_TORQUE}, and the claw only gets 60\% power
(\code{CLAW_SPEED}), because once the claw has closed on an object, extra power
only stalls the motor and makes it overheat. For tuning, \code{CLAW_DIAGNOSTIC}
prints the claw motor's current and temperature to the Brain screen five times
per second. This is wrapped in a \code{try} block (lines 53--58), so that if the
SDK cannot report these values, the diagnostic fails quietly instead of
stopping the drive in the middle of a match.

There are optional software limits on the arm (\code{ARM_USE_LIMITS}). When
enabled, the arm refuses to move further past its lowest or highest position,
but is always allowed to move back.

\section{Cone Assist}
\label{sec:assist}

The cone assist lets the driver press one button to make the robot find the cone
and turn to face it. The left and right arrows start a search to the left or
right. It uses the AI Vision sensor, which is explained in \cref{ch:vision};
here we only describe how the assist behaves.

The assist moves through three states:

\begin{enumerate}
  \item \textbf{Searching:} spin slowly (15\%) in the chosen direction. The speed
    is kept low because the camera's image blurs when the robot spins quickly,
    and a blurred cone is not recognised.
  \item \textbf{Aligning:} once the cone has been seen in 3 frames in a row,
    turn towards it with a power proportional to how far off-centre it is.
  \item \textbf{Centred:} once the cone has been within 4 pixels of the centre for
    3 frames in a row, stop and rumble the controller twice.
\end{enumerate}

If the cone disappears for 5 frames in a row while aligning, the assist goes
back to searching.

\lstinputlisting[firstline=77,firstnumber=77,lastline=82,caption={\code{assist_steer()} in \code{src/drivercontrol.py}, lines 77--82.},label=lst:steer]{\src{drivercontrol.py}}

The turn power is \code{ASSIST_STEER_KP} (0.12) multiplied by the error in
pixels. This is called \emph{proportional control}~\cite{astrom2008}: a big error
gives a fast turn, and a small error gives a slow turn, so the robot slows down
as it approaches the target and does not overshoot. However, a very small power
is not enough to overcome friction, and the robot would stop just short of the
centre. The power is therefore never allowed below 4\% (\code{ASSIST_MIN_TURN}),
and it is capped at 15\% so the image does not blur.

\textbf{Worked example:} If the cone is 100 pixels right of centre, the power is
$0.12 \times 100 = 12\%$. At 10 pixels, $0.12 \times 10 = 1.2\%$, which is raised to
the minimum of 4\%.

\subsection{Cancelling}

The driver must always be able to take back control instantly, because a stuck
assist during a match could cost the game. Every loop, the assist checks all
twelve buttons and both sticks; any button, or any stick pushed past 20, cancels
it straight away.

This created a problem. The button that \emph{starts} the assist is still being
held down on the assist's first loop, so it would immediately cancel the assist.
The \code{released} flag fixes this: the start button is ignored until it has
been let go once. The reverse problem happens in \code{user_control()}: a
button held down after cancelling would immediately start the assist again. The
start is therefore \emph{edge-triggered}, meaning it only fires at the moment the
button goes from not-pressed to pressed:

\lstinputlisting[firstline=174,firstnumber=174,lastline=188,caption={Edge-triggered start in \code{user_control()}, \code{src/drivercontrol.py}, lines 174--188.},label=lst:edge]{\src{drivercontrol.py}}

The assist also uses \code{spin()}, which sets a speed and returns immediately,
rather than \code{spin_for()}, which waits until a movement finishes. A function
that waits cannot check the buttons in the meantime, so it could not be
cancelled.

\section{Summary}

\begin{enumerate}
  \item Stick readings are deadbanded, squared and scaled by a speed mode
    (\sref{sec:stick}).
  \item Slew limiting ramps power up gently but lets it drop quickly
    (\sref{sec:slew}).
  \item The turn is tapered at speed and shrunk before it can clip, keeping
    forward speed (\sref{sec:taper}).
  \item The cone assist searches, aligns with proportional control, and can be
    cancelled by any input (\sref{sec:assist}).
\end{enumerate}

We now explain how the robot recognises the cone.
